% Generated shared reviewed material.
% Full proofs of accepted coherent-band research, with no new target claim.
\section{Proofs for the coherent-band comparisons}
\label{cbk:details}
These proofs use de Faveri's pinned
\href{https://arxiv.org/abs/2610.04045v1}{arXiv:2610.04045v1},
Theorem~1.1 and its stated character family as external input. The
primitive Poisson convention agrees with its Lemma~3.1; the formulas
below also derive the normalization. This is not a proof of that deep
large-sieve theorem. The source version and independent proof reviews
are bound by the accompanying source map; full third-party reading
caches are not public source dependencies.

\subsection{Finite Fourier transforms and the primitive Gaussian}
\label{cbk:gaussian-proof}
For $\mathfrak f=(\gamma)$ use the primitive additive pairing
\[
 e_\gamma(x)=\exp\bigl(2\pi i\operatorname{Tr}(x/(\delta\gamma))\bigr),
 \qquad \delta=1+2\omega.
\]
Let $\theta$ be primitive modulo $\mathfrak f$, extended by zero at
nonunits of this \emph{actual} conductor. For every $\beta\in\mathcal O$,
\begin{equation}\label{cbk:finite-Fourier}
 \sum_{x\bmod\gamma}\theta(x)e_\gamma(\beta x)
          =\overline{\theta(\beta)}\tau_\gamma(\theta).
\end{equation}
At a unit this is change of variable. At a nonunit frequency the
additive character factors through a proper conductor quotient.
Primitivity supplies a unit in the kernel of unit reduction on which
$\theta$ is nontrivial. Multiplication by that unit leaves the additive
factor unchanged and forces the sum to vanish. Unit reduction is
surjective prime by prime and hence by CRT. Parseval for this
nondegenerate finite additive pairing gives
\[
 \sum_{\beta\bmod\mathfrak f}
   \left|\sum_x\theta(x)e_\gamma(\beta x)\right|^2
       =N\mathfrak f\,\varphi(\mathfrak f),\qquad
 |\tau_\gamma(\theta)|^2=N\mathfrak f.
\]
For the unit conductor use its one-element convention $\tau=1$.
For coprime $\gamma_1,\gamma_2$, the two CRT frequencies are
$\gamma_2^{-1}$ and $\gamma_1^{-1}$, so
\begin{equation}\label{cbk:Gauss-CRT}
 \tau_{\gamma_1\gamma_2}(\theta_1\theta_2)
 =\theta_1(\gamma_2)\theta_2(\gamma_1)
                        \tau_{\gamma_1}(\theta_1)\tau_{\gamma_2}(\theta_2).
\end{equation}
The product is primitive because each ramified local conductor occurs
in only one factor.

Here are the archimedean factors that a raw Gauss sum alone would miss.
Put $d=|k|$ and
\[
 P_k(z)=\begin{cases}z^k,&k\ge0,\\\bar z^{-k},&k<0.\end{cases}
\]
With Euclidean Fourier transform
$\widehat g(w)=\int_{\mathbb C}g(z)e^{-2\pi i\operatorname{Re}(z\bar w)},dz$,
differentiating the two-dimensional Gaussian gives
\begin{equation}\label{cbk:harmonic-Fourier}
 \widehat{P_k(z)e^{-\pi b|z|^2}}(w)
    =(-i)^d b^{-d-1}P_k(w)e^{-\pi|w|^2/b},\qquad b>0.
\end{equation}
Holomorphic or antiholomorphic harmonicity removes the lower
contraction terms. The sign is $(-i)^{|k|}$ for either sign of $k$.
Let $\chi_0=\chi N^{-it}$ and
$a_\gamma=2/(\sqrt3|\gamma|)$. Its theta series is
\[
 \Theta_{\chi_0}(t)=\sum_{\alpha\in\mathcal O}
          \theta(\alpha)P_k(\alpha)e^{-\pi a_\gamma t|\alpha|^2}.
\]
The lattice $\gamma\mathcal O$ has covolume
$(\sqrt3/2)|\gamma|^2$ and dual lattice
$2/(\bar\gamma\bar\delta)\mathcal O$.
Write a dual point as $w=2\bar\beta/(\bar\gamma\bar\delta)=c\bar\beta$.
Then $|c|=a_\gamma$ and
$\operatorname{phase}(c)=\operatorname{phase}(\delta\gamma)$.
Equations~\eqref{cbk:finite-Fourier} and
\eqref{cbk:harmonic-Fourier}, applied in each residue class, yield
\begin{equation}\label{cbk:theta-inversion}
 \Theta_{\chi_0}(t)=i^{-d}\operatorname{phase}(\delta\gamma)^k
      \frac{\tau_\gamma(\theta)}{\sqrt{N\mathfrak f}}
                 t^{-d-1}\Theta_{\bar\chi_0}(1/t).
\end{equation}
The inverse different and the conductor generator have not been
replaced by a raw character value at $\delta$.

Use $\Gamma_{\mathbb C}(s)=2(2\pi)^{-s}\Gamma(s)$ and
\[
 \Lambda(s,\chi)=(3N\mathfrak f)^{s/2}
           \Gamma_{\mathbb C}(s-it+d/2)L(s,\chi).
\]
Every nonzero ideal has six generators; their contributions agree
because $\theta(u)u^k=1$. The Mellin transform of the theta series,
with exponent $s+d/2$, is therefore exactly
$3(3N\mathfrak f)^{d/4}\Lambda(s,\chi_0)$.
The same constant occurs on the conjugate side. Mellin integration
of~\eqref{cbk:theta-inversion} proves the root formula for $t=0$.
At nonunit conductor the zero frequency vanishes by primitivity.
At unit conductor and $k\ne0$, $P_k(0)=0$. At unit conductor and
$k=0$, retain the Gaussian zero term when splitting its Mellin
integral at one; this gives the usual meromorphic polar terms and
root one. They are not dropped from the functional equation.

Changing $\gamma$ to $u\gamma$ multiplies its Gauss sum by $\theta(u)$
and its angular factor by $u^k$, whose product is one. Thus the
formula is independent of the conductor generator. Complex
conjugation gives
$\tau_\gamma(\bar\theta)=\theta(-1)\overline{\tau_\gamma(\theta)}$,
where $\theta(-1)=(-1)^k$; together with the angular factor this
proves $\epsilon(\bar\chi)=\overline{\epsilon(\chi)}$.
Finally
\[
 \Lambda(s,\chi_0N^{it})=(3N\mathfrak f)^{it/2}
                                  \Lambda(s-it,\chi_0),\qquad
 \epsilon(\chi_0N^{it})=(3N\mathfrak f)^{it}\epsilon(\chi_0).
\]
This proves~\eqref{cbk:root-formula} also for radial unitary twists.

In~\eqref{cbk:Gauss-CRT}, the angular cross factors
$\operatorname{phase}(\gamma_2)^{k_1}
 \operatorname{phase}(\gamma_1)^{k_2}$ convert the finite residue
values into ideal values $\chi_1(\mathfrak f_2)\chi_2(\mathfrak f_1)$.
The remaining $i$ factor is
$i^{|k_1|+|k_2|-|k_1+k_2|}$. Radial cross factors satisfy
\[
 \prod_j(3N\mathfrak f_j)^{it_j}
       \prod_{i\ne j}(N\mathfrak f_j)^{it_i}
       =\left(3N\prod_j\mathfrak f_j\right)^{i\sum_jt_j}.
\]
Iteration proves~\eqref{cbk:root-product}. Coprimality is a hypothesis
on the \emph{actual primitive ideal conductors}, not on rational
norms or higher induced moduli.

\subsection{Transport to every original class and ray component}
\label{cbk:root-transport-proof}
Let the types of $\xi_{\mathfrak m},\xi_{\mathfrak n}$ be
$(k_A,t_A),(k_B,t_B)$. The character $\sigma$ in~\eqref{cbk:eta} has
types $(k_B-k_A,t_B-t_A)$. The three primitive conductors
$\mathfrak m,\mathfrak n,\mathfrak f_S$ are pairwise coprime;
no good local ramification can cancel against the other two.
Applying~\eqref{cbk:root-product} gives
\begin{equation}\label{cbk:triple-root}
 \begin{split}
 \epsilon(\eta)
 &=\gamma_{AB\rho}\epsilon(\xi_{\mathfrak m})
       \overline{\epsilon(\xi_{\mathfrak n})}
       \xi_{\mathfrak m}(\mathfrak n\mathfrak f_S)
       \overline{\xi_{\mathfrak n}(\mathfrak m\mathfrak f_S)}
                                  \sigma(\mathfrak m\mathfrak n),\\
 \gamma_{AB\rho}
 &=i^{|k_A|+|k_B|+|k_B-k_A|}\epsilon(\sigma).
 \end{split}
\end{equation}
Source reciprocity is
$\widetilde\chi_{\mathfrak m}(\mathfrak n)
 \overline{\widetilde\chi_{\mathfrak n}(\mathfrak m)}=\alpha(A,B)$.
Substituting $\widetilde\chi=\xi\psi$ in~\eqref{cbk:triple-root}
cancels all crossed $\psi$ values and proves~\eqref{cbk:two-gauges},
with $\kappa_{AB\rho}=\alpha(A,B)\gamma_{AB\rho}$.
In particular the remaining factor is $\rho(\mathfrak m)\rho(\mathfrak n)$,
not an assumed class constant; $\rho$ need not factor through the
source class group.

For good dual $\mathfrak z$, write
$\widetilde\chi_{\mathfrak m}(\mathfrak z)
 =\gamma_A(\mathfrak z)R_{\mathfrak m}(\mathfrak z)$,
with the fixed finite reciprocity phase in that class. Then
\begin{equation}\label{cbk:Omega}
 \begin{aligned}
 \overline{\sigma(\mathfrak z)}
   \overline{\xi_{\mathfrak m}(\mathfrak z)}\xi_{\mathfrak n}(\mathfrak z)
 &=\Omega_{AB\rho}(\mathfrak z)
                  \overline{R_{\mathfrak m}(\mathfrak z)}
                            R_{\mathfrak n}(\mathfrak z),\\
 \Omega_{AB\rho}
 &=\bar\sigma\,\psi_A\bar\psi_B\,\bar\gamma_A\gamma_B,
 \qquad |\Omega_{AB\rho}(\mathfrak z)|=1.
 \end{aligned}
\end{equation}
The equality includes raw collision zeros. This is a fixed unitary
\emph{row} phase for each component, not equality of the two column
classes. At $S$ the separate scalar
$\sigma(\mathfrak k^2\mathfrak l\mathfrak e)
 \overline{\sigma(\mathfrak h)}$ is kept, including zero.
There is no division by such a scalar, and no substitution of
$\rho(\mathfrak e/\mathfrak h)$ at a conductor drop.
Off the original coprime column support, products of individual
column gauges below are only an auxiliary extension. They are not
roots or primitive conductors of overlapping pairs.

\subsection{The sharp capacity, square-divisor signs and exact Poisson}
\label{cbk:kernel-proof}
The endpoints in~\eqref{cbk:nodal} are integers. The two smooth
transitions lie between ideal norm nodes, so its value at
$x=N\mathfrak v_1/A_1$ is exactly one or zero according to the
original interval and capacity. For nonempty fibres these transitions
are disjoint. Their monotonicity gives total variation two; their
$j$th derivatives cost powers of $A_1$.
The elementary identity
\begin{equation}\label{cbk:row-inversion}
 \sum_{\mathfrak v\ \mathrm{good}}\mu(\mathfrak v)^2
                    \mathbf1_{(\mathfrak v,\mathfrak v_2)=1}B(\mathfrak v)
 =\sum_{\substack{\mathfrak k\ \mathrm{good,SF}\\
                    (\mathfrak k,\mathfrak v_2)=1}}\mu(\mathfrak k)
   \sum_{\mathfrak l\mid\mathfrak v_2}\mu(\mathfrak l)
       \sum_{\mathfrak z\ \mathrm{good}}B(\mathfrak k^2\mathfrak l\mathfrak z)
\end{equation}
follows by expanding the squarefree indicator and then the remaining
coprimality indicator. Auxiliary $\mathfrak z$ need not be squarefree,
and may overlap $\mathfrak k$ or $\mathfrak l$. The original norm
support fixes $N\mathfrak l\le b$ and
$N\mathfrak k\le\sqrt{b/N\mathfrak l}$ \emph{before} Poisson.
Including more $\mathfrak k$ would add zero primal terms whose
individual dual windows need not be zero.

On good $\mathfrak v_1$, expand exactly
\[
 \mathbf1_{\rm ray}R_{\mathfrak m}(\mathfrak v_1)
       \overline{R_{\mathfrak n}(\mathfrak v_1)}
   =\sum_\rho h_{AB\rho}(\mathfrak v_2)\eta_{\mathfrak m,\mathfrak n,\rho}
                                                    (\mathfrak v_1).
\]
This includes every reciprocity/ray component, not a selected same
class term. With $L=A_1/((N\mathfrak k)^2N\mathfrak l)$, good-$S$
deletion is the \emph{finite} identity
\begin{equation}\label{cbk:S-deletion}
 \sum_{\mathfrak z\ \mathrm{good}}G(N\mathfrak z/L)\eta(\mathfrak z)
 =\sum_{\mathfrak e\mid P_S}\mu(\mathfrak e)\eta(\mathfrak e)
       \sum_{\mathfrak y}G(N\mathfrak y/(L/N\mathfrak e))\eta(\mathfrak y).
\end{equation}
Primitive zeros at $S$ suppress the corresponding terms automatically.
They do not assign forbidden raw values to an original row.

For a nonprincipal primitive type-zero $\eta$ of conductor norm $F$,
its normalized functional equation and ordinary Mellin inversion give
\begin{equation}\label{cbk:primitive-Poisson}
 \sum_{\mathfrak y}G(N\mathfrak y/L)\eta(\mathfrak y)
 =\epsilon(\eta)\frac L{\sqrt F}
      \sum_{\mathfrak y}\dot G(LN\mathfrak y/F)
                                      \overline{\eta(\mathfrak y)}.
\end{equation}
To see the field constant explicitly, put
\[
 \gamma_K(s)=C_0^{s-1/2}\frac{\Gamma(s)}{\Gamma(1-s)},\qquad
 \widehat G(s)=\int_0^\infty G(t)t^s\,\frac{dt}{t}.
\]
The Mellin transform of~\eqref{cbk:bessel} is
$\gamma_K(s)\widehat G(1-s)$. The Bessel identity may first be
justified absolutely in $0<\operatorname{Re}s<1/4$; smooth compact
support and repeated Hankel integration by parts continue it to
$\operatorname{Re}s>0$ at each fixed $G$. Mellin inversion of the
primal ideal sum, its primitive functional equation, and the change
$s\mapsto1-s$ now give~\eqref{cbk:primitive-Poisson} with exactly
$L/\sqrt F$. There is no principal residue since $\eta$ has nonempty
good ramification. The intermediate contour moves involve the
entire nonprincipal completed function and the fixed smooth
transform, not an inverse $L$ function or an assumed zero-free strip.
This also recovers the pinned source Lemma~3.1 convention.

At a good prime ideal of norm $P$, a nontrivial primitive character
modulo $p$ lifted to $p^j$, $j\ge2$, has raw Gauss sum zero for a
primitive additive frequency. More generally it vanishes unless the
frequency has valuation $j-1$, when its value is
$P^{j-1}\bar\chi(\beta/\pi^{j-1})\tau_p(\chi)$ with compatible
additive normalizations. This follows by summing the lifts first.
Thus neither lower character order nor a higher physical modulus
licenses a unit root. Applying~\eqref{cbk:primitive-Poisson} only at
the actual $F$ in~\eqref{cbk:eta}, and inserting
\eqref{cbk:S-deletion} into~\eqref{cbk:row-inversion}, proves the exact
signed kernel. Splitting every dual ideal uniquely as
$\mathfrak y=\mathfrak h\mathfrak z$, with $\mathfrak h$ supported
on $S$ and $\mathfrak z$ good, retains all modes.

\subsection{The joint bilinear estimate behind finite prefixes}
\label{cbk:prefix-proof}
The external sixth-order norm gives the following arbitrary-column-
coefficient comparison, for every $Z,L\ge1$:
\begin{equation}\label{cbk:all-ideal-norm}
 \sum_{\substack{\mathfrak z\ \mathrm{good}\\N\mathfrak z\le Z}}
    \left|\sum_{\substack{\mathfrak a\ \mathrm{SF,good}\\N\mathfrak a\asymp L}}
                      b_{\mathfrak a}R_{\mathfrak a}(\mathfrak z)\right|^2
       \ll (ZL)^\epsilon\Phi(Z,L)\sum_{\mathfrak a}|b_{\mathfrak a}|^2.
\end{equation}
For completeness, extract $\mathfrak z=\mathfrak r\mathfrak s^6$
with sixth-power-free $\mathfrak r$. Freeze $\mathfrak s$ as the
original column collision mask $(\mathfrak a,\mathfrak s)=1$ and
apply~\eqref{phy:input} to $\mathfrak r$. The four sums have factors
$\sum(N\mathfrak s)^{-6}$, the count
$N\mathfrak s\le Z^{1/6}$, $\sum(N\mathfrak s)^{-5}$ and
$\sum(N\mathfrak s)^{-2}$, respectively. Ideal counting proves
\eqref{cbk:all-ideal-norm}. This is not an arbitrary \emph{row}-weight
theorem or an extension of a residue-subtracted energy estimate.
For $1/2\le L<1$ use the direct unit residual column; below $1/2$
there is no annular residual column. Primitive $\xi$ is not treated
as sixth-order: convert to raw $R$ inside each fixed residual-class
positive norm, where its common $\psi$ row phase cancels.

At $s=\lambda+i\tau$, use~\eqref{cbk:two-gauges} to form the two
\emph{distinct} vectors
\begin{equation}\label{cbk:Mellin-vectors}
 \begin{aligned}
 b_{\mathfrak m}(s)&=c_{\mathfrak m}u_{\mathfrak m}
  R_{\mathfrak m}(\mathfrak v_2^2)
  \xi_{\mathfrak m}(\mathfrak k^2\mathfrak l\mathfrak e)
  \overline{\xi_{\mathfrak m}(\mathfrak h)}(N\mathfrak m)^{s-1/2},\\
 d_{\mathfrak n}(s)&=c_{\mathfrak n}v_{\mathfrak n}
  R_{\mathfrak n}(\mathfrak v_2^2)
  \xi_{\mathfrak n}(\mathfrak k^2\mathfrak l\mathfrak e)
  \overline{\xi_{\mathfrak n}(\mathfrak h)}(N\mathfrak n)^{s-1/2}.
 \end{aligned}
\end{equation}
The pairing is $b_{\mathfrak m}(s)
\overline{d_{\mathfrak n}(\bar s)}$, not a PSD form at $s$.
The exact remaining primitive-conductor factor is
\begin{equation}\label{cbk:FS-Mellin}
 (N\mathfrak f_S)^{s-1/2}
       \left(\frac{A_1}{N\mathfrak l N\mathfrak e}\right)^{1-s},
\end{equation}
with scalar $\sigma(\mathfrak k^2\mathfrak l\mathfrak e)
\bar\sigma(\mathfrak h)$ and finite $\mathfrak k$ weight
$(N\mathfrak k)^{-2+2s}$. Its $S$ exponent is $s-1/2$, not its
inverse. On $0<\lambda<1/2$ its modulus is at most one.

Insert exact gcd inversion before comparison and put
$\mathfrak m=\mathfrak d\mathfrak a$,
$\mathfrak n=\mathfrak d\mathfrak b$.
Residual $\mathfrak a,\mathfrak b$ may overlap each other but are
individually coprime to $\mathfrak d$. Raw multiplicativity gives
\[
 \overline{R_{\mathfrak d\mathfrak a}(\mathfrak z)}
 R_{\mathfrak d\mathfrak b}(\mathfrak z)
 =\mathbf1_{(\mathfrak d,\mathfrak z)=1}
              \overline{R_{\mathfrak a}(\mathfrak z)}R_{\mathfrak b}(\mathfrak z).
\]
Their source classes are $A[\mathfrak d]^{-1}$ and
$B[\mathfrak d]^{-1}$, generally different. Bilinear Cauchy in the
row variable, including~\eqref{cbk:Omega}, gives two separate
positive norms. Only \emph{inside} them may the raw
$(\mathfrak d,\mathfrak z)$ mask be dropped. Uniformly in the
fixed column phases and $\tau$,
\begin{equation}\label{cbk:subsequence}
 \sum_{\mathfrak a}|b_{\mathfrak d\mathfrak a}(s)|^2,
 \quad\sum_{\mathfrak b}|d_{\mathfrak d\mathfrak b}(\bar s)|^2
       \ll B_W^2D^{2\lambda+\epsilon}/N\mathfrak d.
\end{equation}
Indeed $|A_U(\mathfrak m)|\le3^{\omega(\mathfrak m)}$ and an annular
subsequence has $O(D/N\mathfrak d)$ columns. Neither
$A_U(\mathfrak d\mathfrak a)$ nor $\beta_{\mathfrak d\mathfrak a}$
is factored. For the small residual interval the direct unit count
has the same bound.

Apply~\eqref{cbk:all-ideal-norm} to the two separate vectors. The
divisor sums have exponents $1,2,4/3,11/6$:
\begin{equation}\label{cbk:d-Phi}
 \sum_{N\mathfrak d\le2D}\frac1{N\mathfrak d}
                      \Phi(Z,D/N\mathfrak d)
      \ll\log(2D)\Phi(Z,D),
\end{equation}
with the direct endpoint understood. For
$Y<N\mathfrak h N\mathfrak z\le2Y$, the weight modulus is at most
$Y^{-\lambda}$ and the legal positive row prefix is
$N\mathfrak z\le2Y/N\mathfrak h\ge1$.
The four positive row exponents in $\Phi$ give
\[
 \sum_{\mathfrak h\ S\text{-supported},\,N\mathfrak h\le2Y}
                  \Phi(2Y/N\mathfrak h,D)\ll_S\Phi(Y,D).
\]
This is a summation over \emph{all} $S$-supported modes, not an
exponent-zero convergent sum or a missing raw $S$ value.

The sharp $G$ requires an honest transform bound. Integration by
parts against its finite variation gives
\[
 \widehat G(1-s)=-\frac1{1-s}\int t^{1-s}\,dG(t).
\]
For fixed $0<\lambda<1/2$, Stirling's ratio estimate is
$|\gamma_K(\lambda+i\tau)|\ll_\lambda
(1+|\tau|)^{2\lambda-1}$. Thus
\begin{equation}\label{cbk:BV-L1}
 \int_{\mathbb R}|\widehat G(1-\lambda-i\tau)
                          \gamma_K(\lambda+i\tau)|\,d\tau
                         \ll_\lambda1,
\end{equation}
using the integrable tail $(1+|\tau|)^{2\lambda-2}$ and
$\operatorname{TV}(G)\le2$. No derivative seminorm growing with
$A_1$ is hidden here.
Only after the joint estimate may the original finite $\mathfrak k$
sum be bounded by $\zeta_K^S(2-2\lambda)$; the identity has not been
Euler-completed. Finite $\mathfrak e$, the original
$\mathfrak l\mid\mathfrak v_2$ divisor sum, and unweighted actual
$\mathfrak v_2$ counting cost a fixed constant and
$A_2(DA_2)^\epsilon$. Equations~\eqref{cbk:subsequence}--
\eqref{cbk:BV-L1} give
\begin{equation}\label{cbk:fixed-line-band}
 \frac{|\mathcal P_D(Y,2Y]|}{D}
 \ll_{\lambda}B_W^2(DYA_2)^\epsilon
    A_2A_1^{1-\lambda}D^{2\lambda-1}Y^{-\lambda}\Phi(Y,D).
\end{equation}
Every finite source/ray component is reconstructed using the distinct
gauges and bilinear argument, rather than a same-class inference.

For $0<\lambda<1/6$, each of the four powers of $Y$ after
multiplication by $Y^{-\lambda}$ is positive. Finite dyadic summation,
a last partial band and the dual unit separately therefore give
the same bound with $Y=X$ for every finite prefix $X\ge1$.
For a requested final loss $\epsilon_0>0$, fix
$\lambda=\min\{1/12,\epsilon_0/4\}$, independently of all scales,
and use the positive-norm loss $\epsilon_0/2$. Since $A_1,X\ge1$,
\[
 D^{2\lambda}(A_1X)^{-\lambda}\le D^{\epsilon_0/2},\qquad
 (DXA_2)^{\epsilon_0/2}D^{\epsilon_0/2}\le(DXA_2)^{\epsilon_0}.
\]
This proves~\eqref{cbk:prefix-bound}; neither endpoint
$\lambda=0$ nor $\lambda=1/2$ is invoked.

For the independent all-orders bound, expand $J_0$ exactly at every
finite dual prefix. The $j$th moment satisfies
$|M_j|=|\int t^jG(t)dt|\le2(5/2)^j$.
The resulting original-column factors $(D/N\mathfrak m)^j$ and
$(D/N\mathfrak n)^j$ are contractions, including on
$\mathfrak m=\mathfrak d\mathfrak a$; they are not replaced by
$(D/N\mathfrak a)^j$. The exact $S$ factor is
$(N\mathfrak f_S)^{-j-1/2}$, and the finite $\mathfrak k$ norm
exponent is $2+2j$ with the \emph{same} character phase
$\eta(\mathfrak k)^2$ as before. After gcd inversion, weighted
bilinear Cauchy may bound
$(N\mathfrak h N\mathfrak z)^j\le X^j$ inside each positive norm;
the product of the square roots costs $X^j$, not $X^{2j}$.
All remaining divisor/$S$/subsequence constants are uniform in $j$.
The $j$th joint bound is the prefactor of~\eqref{cbk:all-j} times
$((5/2)A_1X/(C_0D^2))^j/(j!)^2$. Summing this absolutely convergent
series proves~\eqref{cbk:all-j} without an infinite-dual interchange.

For completeness, the scoped critical-half estimate uses
$H(v)=e^{v/2}G(e^v)$. Its Fourier transform is
$\widehat G(1/2-i\tau)$. The fixed smooth step and its two
transitions give
$\|H\|_1+\|H'\|_1\ll1$ and $\|H''\|_1\ll1+A_1$.
The minimum of the three Fourier bounds
$O(1)$, $O(1/|\tau|)$ and $O(A_1/\tau^2)$ has integral
$O(1+\log(2A_1))$. On the half line
$|\gamma_K(1/2+i\tau)|=1$ exactly. The original finite square-divisor
sum is harmonic, $\sum_{N\mathfrak k\le K_{\mathfrak l}}
(N\mathfrak k)^{-1}\ll1+\log(2A_1)$; it is not an infinite
$\zeta_K(1)$. The same joint argument in the stated same-class/
$\rho=1$ component proves~\eqref{cbk:half}. This estimate is not
summed over all dual bands to infer a whole target.

\subsection{Finite-height Schur transfer, including the operator error}
\label{cbk:schur-proof}
For $x\ne1$ define
\[
 I_T(x)=\frac1{2\pi}\int_{-T}^T\frac{x^{c_P+it}}{c_P+it}\,dt.
\]
Fourier inversion of $e^{-c_Pv}\mathbf1_{v\ge0}$ gives its full
vertical limit $\mathbf1_{x>1}$. Integration by parts in each
oscillatory tail gives
\begin{equation}\label{cbk:Perron-error}
 |I_T(x)-\mathbf1_{x>1}|
                \le\frac{2x^{c_P}}{\pi T|\log x|}.
\end{equation}
There is no pole-free inverse-$L$ contour assertion. Set
$U^\#=\lfloor U\rfloor+1/2$. Every divisor norm $n\le X$ is a
positive integer, so
\[
 |\log(U^\#/n)|\ge\frac1{2X+1}\ge\frac1{3X},\qquad
 (U^\#/n)^{c_P}\le e.
\]
The indicator is exactly $n\le U$, even when $U$ is an integer.
Since $6e/\pi<6$,~\eqref{cbk:Perron-error} is at most
$e_T=6X/T\le3/4$. A product of the two indicators therefore has
error at most $2e_T+e_T^2\le18X/T$.

For squarefree $\mathfrak m$ put
\[
 H_{\mathfrak m}(v,w)=\prod_{p\mid\mathfrak m}
                              (1-(Np)^{-v}-(Np)^{-w}).
\]
There are exactly $3^{\omega(\mathfrak m)}$ ordered disjoint divisor
assignments. On the supported coprime products only,
$H_{\mathfrak d\mathfrak a}=H_{\mathfrak d}H_{\mathfrak a}$, so
\begin{equation}\label{cbk:finite-Perron}
 \begin{split}
 A_U(\mathfrak d\mathfrak a)
 &=\frac1{(2\pi i)^2}\int_{c_P-iT}^{c_P+iT}
                \int_{c_P-iT}^{c_P+iT}
       \frac{(U^\#)^{v+w}}{vw}
                 H_{\mathfrak d}(v,w)H_{\mathfrak a}(v,w)\,dv\,dw
                           +R_T(\mathfrak d,\mathfrak a),\\
 |R_T(\mathfrak d,\mathfrak a)|
 &\le(18X/T)3^{\omega(\mathfrak d)}3^{\omega(\mathfrak a)}.
 \end{split}
\end{equation}
Forbidden entries of $K$ remain zero. The integral is the actual
operator $\operatorname{diag}(H_{\mathfrak d})K
\operatorname{diag}(H_{\mathfrak a})$, with an ordinary right
diagonal, not a Hermitian assumption. The two diagonal norms cost
$G_{\mathcal I\mathcal J}$, $(U^\#)^{2c_P}\le e^2$, and each
normalized absolute vertical integral is $L_T$.

A uniform entry error alone does not multiply operator norms.
Instead, with counting $\ell^2$ norms,
\[
 \|R_T\circ K\|_{\rm op}
  \le\|R_T\circ K\|_F
  \le\max|R_T|\,\|K\|_F
  \le\max|R_T|\sqrt{\min(\#\mathcal I,\#\mathcal J)}\,\|K\|_{\rm op}.
\]
The accepted fixed-field ideal bound $\#\{N\mathfrak a\le X\}\le2X$
gives an error at most $18\sqrt2X^{3/2}G_{\mathcal I\mathcal J}/T$,
less than the displayed $32$ constant. This proves
\eqref{cbk:schur-bound}. Splitting the finite set of prime ideals
$Np<3^{1/\eta}$ from the rest proves
$3^{\omega(\mathfrak m)}\ll_\eta(N\mathfrak m)^\eta$.
Use $\eta=\epsilon/2$ and $T=(4X)^4$ to obtain the subpower/log bound.
This is a mathematical finite-matrix transfer, not a claim of
fast numerical Perron evaluation at that height.

Let $F_0(v)=W(e^v)$. Two integrations by parts give
$|\widehat F_0(t)|\le\min\{\|F_0\|_1,\|F_0''\|_1/t^2\}$.
Split at $|t|=1$ to get
\[
 M(W)\le(\|F_0\|_1+\|F_0''\|_1)/\pi.
\]
Here $F_0''(v)=xW'(x)+x^2W''(x)$ and its support has length
$\log2$, proving~\eqref{cbk:profile-seminorm}.
Absolute Fourier inversion expresses $W(N\mathfrak dN\mathfrak a/D)$
as an integral of two unitary norm diagonals and a scalar. Its
Schur multiplier norm is at most $M(W)$. Compose with
\eqref{cbk:finite-Perron}; two coefficient occurrences cost the
square. All original zeros remain throughout, and no factorization
of a primitive $\beta_{\mathfrak d\mathfrak a}$ is supplied by this lemma.

\subsection{The cubic transposition in an identity divisor class}
\label{cbk:cubic-proof}
Fix the scope of Proposition~\ref{cbk:cubic}. On good arguments the
primary-generator source construction is index-multiplicative.
For $[\mathfrak d]=1$, comparison with the unit source class removes
the $S$ conductor of $\widetilde\chi_{\mathfrak d}$; its type is zero.
For $[\mathfrak a]=C$ and $(\mathfrak d,\mathfrak a)=1$,
\[
 \xi_{\mathfrak d\mathfrak a}
       =\widetilde\chi_{\mathfrak d}\xi_{\mathfrak a}.
\]
The primitive conductors are coprime and their angular correction in
\eqref{cbk:root-product} is one. Multiplication by the actual
$\psi_C(\mathfrak d\mathfrak a)$ yields
\[
 \beta_{\mathfrak d\mathfrak a}
 =\epsilon(\widetilde\chi_{\mathfrak d})\beta_{\mathfrak a}
        \widetilde\chi_{\mathfrak d}(\mathfrak a)
        \widetilde\chi_{\mathfrak a}(\mathfrak d).
\]
Source reciprocity is
$\widetilde\chi_{\mathfrak d}(\mathfrak a)
\overline{\widetilde\chi_{\mathfrak a}(\mathfrak d)}=\alpha(1,C)$,
proving~\eqref{cbk:beta-split}. In the product for residuals
$\mathfrak a,\mathfrak b$, their common $\epsilon$ and $\alpha$ cancel:
\begin{equation}\label{cbk:cubic-quotient}
 \beta_{\mathfrak d\mathfrak a}
       \overline{\beta_{\mathfrak d\mathfrak b}}
 =\beta_{\mathfrak a}\overline{\beta_{\mathfrak b}}
       (\widetilde\chi_{\mathfrak a}
                    \overline{\widetilde\chi_{\mathfrak b}})^2(\mathfrak d).
\end{equation}
If $\mathfrak g=(\mathfrak a,\mathfrak b)$, the primitive cubic
conductor is $(\mathfrak a/\mathfrak g)(\mathfrak b/\mathfrak g)$,
nonprincipal when $\mathfrak a\ne\mathfrak b$. The same source class
cancels its $S$ component. Common good primes have dropped only from
that primitive conductor; the full $\mathfrak a\mathfrak b$ union
still remains excluded from $\mathfrak d$.

The other common $\mathfrak d$ phases from $\nu,u_0$ and the
$\mathfrak k,\mathfrak l,\mathfrak v_2,\mathfrak e,\mathfrak h$
modulations cancel in the pair. Their raw zeros, together with the
$\mathfrak z$ factor, retain
\[
 (\mathfrak d,S\mathfrak a\mathfrak b\mathfrak k
                                      \mathfrak v_2\mathfrak z)=1.
\]
This notation means coprimality with the union of the displayed
prime supports. Primitive values at $\mathfrak e,\mathfrak h$ are
legitimate unitary phases of a good-conductor character.
The radial factor $(N\mathfrak d)^{2s-1}$ is kept, not called a
cancelled phase.

At fixed $\mathfrak z,s$ the core matrix is
\begin{equation}\label{cbk:cubic-matrix}
 K^{(\mathfrak z)}_{\mathfrak d\mathfrak a}
 =\mathbf1_{[\mathfrak d]=1,[\mathfrak a]=C,
                     (\mathfrak d,\mathfrak k\mathfrak v_2\mathfrak z)=1}
      \widetilde\chi_{\mathfrak a}(\mathfrak d)^2
                                \overline{R_{\mathfrak a}(\mathfrak z)}.
\end{equation}
Its cubic character already has the $(\mathfrak d,\mathfrak a)$ zero.
All additional fixed residual phases are column contractions. On
good arguments the squared sextic symbol is the actual primary
cubic symbol $(\alpha_{\mathfrak a}/\mathfrak d)_3$.
The pinned external Example~2.2 constructs precisely this family
with primary generators congruent to one modulo $3$. Both variables
are squarefree and hence cube-free; source-sixth-class selectors
and deleting the prime above $2$ only restrict the positive matrix.
Theorem~1.1 at $n=3$, or its transpose, therefore gives
\begin{equation}\label{cbk:cubic-norm}
 \|K^{(\mathfrak z)}\|_{\rm op}^2
                \ll_\eta(EL)^\eta\Theta_3(E,L).
\end{equation}
The lengths are physical ideal norms $E,L$, not varying conductor
labels.

The original annular condition implies
$L/2\le N\mathfrak a\le2L$, a bounded number of dyads with lower
endpoint at least one. On the encompassing rectangle
$N\mathfrak dN\mathfrak a\le4D$, apply
Proposition~\ref{cbk:schur} with $X=4D$. This is the necessary
repair: $A_U(\mathfrak d\mathfrak a)$ depends on \emph{both} matrix
indices and cannot be put into a fixed coefficient vector by fiat.
Composing the common $W$ costs $M(W)$.
With $0<\lambda<1/2$, the squared row radial factor is
$O(E^{2\lambda-1})$, and the column squared norm is
\[
 \sum_{L/2\le N\mathfrak a\le2L}(N\mathfrak a)^{2\lambda-1}
                                      \ll_\lambda L^{2\lambda}.
\]
The product is $D^{2\lambda}/E$. Thus the two Mellin vectors have
squared norms at most
\[
 M(W)^2D^{2\lambda+\eta}\Theta_3(E,L)/E.
\]
Cauchy in $\mathfrak d$ now uses $|\mu(\mathfrak d)|\le1$, explicitly
losing rather than exploiting its sign. A final triangle over
$N\mathfrak z\le Z$ gives the new inner branch
\begin{equation}\label{cbk:cubic-inner}
 |T_E(Z;s)|\ll M(W)^2D^{2\lambda+\eta}
                                      (Z/E)\Theta_3(E,L).
\end{equation}
The accepted separate positive-divisor argument gives
$|T_E(Z;s)|\ll M(W)^2D^{2\lambda+\eta}\Phi(Z,L)$: there are
$O(E)$ divisors and each actual subsequence norm is
$O(M(W)^2D^{2\lambda+\eta}/E)$. Both estimates apply to the same
quantity, so their minimum is legal.

For $Y<N\mathfrak hN\mathfrak z\le2Y$, put
$Z_{\mathfrak h}=2Y/N\mathfrak h\ge1$.
The new branch uses the convergent positive exponent-one sum
$\sum_{\mathfrak h}Z_{\mathfrak h}\ll_SY$; the old branch uses
its four positive $\Phi$ exponents. Sum each branch to obtain the
minimum after summation. The same finite $\mathfrak k$ majorant,
$\mathfrak e/\mathfrak l$ bounds, actual $A_2$ count and fixed-line
BV integral as above now prove~\eqref{cbk:cubic-bound}, including
the once-only $1/D$ normalization. No all-divisor or whole-dual
conclusion follows.

On rough prime divisors $\mathfrak d=p$, $Np>U$, no short divisor
can contain $p$, so $A_U(p\mathfrak a)=A_U(\mathfrak a)$ exactly.
The remaining fixed-residual arithmetic sum has the shape
\begin{equation}\label{cbk:prime-failure}
 \begin{split}
 -A_U(\mathfrak a)A_U(\mathfrak b)
 \sum_{\substack{E<Np\le2E,\ Np>U,\ [p]=1\\
      p\notin S\cup\operatorname{supp}(\mathfrak a\mathfrak b
                    \mathfrak k\mathfrak v_2\mathfrak z)}}
 &\lambda_{\mathfrak a\mathfrak b}(p)(Np)^{2s-1}\\[-2mm]
 &\cdot W(NpN\mathfrak a/D)\overline{W(NpN\mathfrak b/D)},
 \end{split}
\end{equation}
up to the already fixed residual phases, with
$\lambda_{\mathfrak a\mathfrak b}
 =(\widetilde\chi_{\mathfrak a}\overline{\widetilde\chi_{\mathfrak b}})^2$.
Here $\mu(p)=-1$ is constant. A fixed-field scalar Mertens estimate
is not a uniform bound for these varying cubic prime twists.
The auxiliary principal case $\mathfrak a=\mathfrak b$ may cancel
against other divisor sectors; it is no physical lower obstruction.

\subsection{Classification and proof of the full equal-norm sector}
\label{cbk:equal-proof}
Factor the common rational norm prime by prime. At an inert prime
there is only one ideal of norm $p^2$; equal norms force the same
squarefree valuation in both columns and coprimality forces it to
zero. The same argument excludes a ramified prime. At a split prime
write the two pairs of valuations as $(x,y),(z,w)\in\{0,1\}^2$.
Then $x+y=z+w$ and $xz=yw=0$. Sum two is forbidden; sum one forces
opposite choices. Hence $\mathfrak n=\bar{\mathfrak m}$ and the
common norm is rational-squarefree supported on split primes.
Both primes above each such prime must already be good. This does
not assume that the original $S$ is conjugation-stable or enlarge it.
Each eligible $\mathfrak m$ has exactly one partner, so the ordered
pair count is at most $\#\{N\mathfrak m\le2D\}\le4D$.
All other rational-nonsquarefree norms remain in the \emph{full}
annulus; only this sector has been classified.

The ternary bound is uniform in the original cutoff:
$|A_U(\mathfrak m)|\le3^{\omega(\mathfrak m)}\ll_\eta D^\eta$.
For each actual row the absolute triangle in this sparse sector
costs $O(B_W^2D^{1+\epsilon})$ before normalization. Divide once
by $D$ and sum the actual rows to obtain~\eqref{cbk:equal-target}.
Original row masks give $R_w\ll R$, and a nonempty capacity block
has $Nu_0A_1A_2^2\le H$, hence $R\le H/(Nu_0A_2)\le H$.
For each fixed $W$, $S_1^{AA}=S_{\rm eq}+S_{\rm neq}$ proves the
target equivalence stated in the core, with constants allowed to
depend on that fixed profile.

Conjugating the residue-field definition gives, including zeros,
\begin{equation}\label{cbk:norm-pullback}
 R_{\bar{\mathfrak m}}(x)=\overline{R_{\mathfrak m}(\bar x)},\qquad
 R_{\mathfrak m}(x)\overline{R_{\bar{\mathfrak m}}(x)}
                                  =R_{\mathfrak m}(Nx).
\end{equation}
For a split rational $p$, $p\mid Nx$ exactly when $x$ meets at
least one of its two conjugate primes, so the right side retains
the original \emph{union} zero. For an orientation of rational-
squarefree norm $N$, restriction to rational integers gives a
primitive Dirichlet character $\psi_{\mathfrak m}\pmod N$:
$\mathbb Z/N\mathbb Z\longrightarrow\mathcal O/\mathfrak m$ is an
isomorphism, and each local character has exact order six.
Its ideal norm pullback has conductor
$(N)=\mathfrak m\bar{\mathfrak m}$ and infinity type zero.
There is no smaller radical or common truncated-frequency assertion.

Choose one prime $\pi_p$ above each admissible split rational $p$.
Every divisor of an orientation $\mathfrak m$ has a unique rational
divisor norm and the same M\"obius parity. Thus
$A_U(\mathfrak m)=a_U(N)$ with the \emph{two original} cutoffs, and
the partner has the same coefficient. At
$t=Nu_0N\mathfrak v_1(N\mathfrak v_2)^2$,
\eqref{cbk:norm-pullback} makes the local orientation factors
$a_p\chi_p(t)$ and its conjugate. Expanding their product gives
all $2^{\omega(N)}$ ordered orientations exactly, proving
\eqref{cbk:orientation}. Both reversals are included, not divided
by two. The common complex profile gives $|W(N/D)|^2$.
For $\nu=1$ and no raw collision the character factors are
$2\operatorname{Re}\mu_6\in\{2,1,-1,-2\}$, so they need not
vanish or be positive; $a_U$ or $W$ can still vanish. This is not
a row-energy lower bound.

For~\eqref{cbk:equal-prefix}, use the exact primitive representation
\emph{restricted first} to these $O(D)$ pairs. Its conductor norm
is $N\mathfrak f_S N^2\ge D^2$, and $|\dot G|\le2C_0^{-1/2}$
by $|J_0(t)|\le1$ on real $t$. The Poisson amplitude is at most
$A_1/((N\mathfrak k)^2N\mathfrak lN\mathfrak eD)$.
There are $O(Y)$ integral dual ideals in a prefix or band, counting
every $S$-supported ideal once. The finite square-divisor sum is
bounded by its convergent positive $Nk^{-2}$ majorant; finite
$S$/ray and $\mathfrak l$ divisor costs give
$A_2(DA_2)^\epsilon$. With the sparse pair count and coefficient
product $B_W^2D^\epsilon$, the numerator is
$O(B_W^2(DYA_2)^\epsilon RY)$. Division by $D$ proves the bound.
This uses no general row-weight norm and no root cancellation.

At fixed parameters, exact Poisson gives
$D S_{\rm eq}=\mathcal P_{\rm eq}(\le Y)+\mathcal P_{\rm eq}(>Y)$,
with its actual smooth derivative constants. Therefore
\begin{equation}\label{cbk:equal-complement}
 \frac{|\mathcal P_{\rm eq}(>Y)|}{D}
 \ll B_W^2\{D^\epsilon R_w+(DYA_2)^\epsilon RY/D\}.
\end{equation}
This is a \emph{signed} whole-minus-prefix identity and triangle,
not absolute summation of high modes. The condition $RY/D\le H$
gives pure $Y\asymp D$ and mixed $Y\asymp D^{7/6}$ target-sized
prefix/complement sectors. Unequal norms and arbitrary larger
individual windows remain unestimated.

\subsection{A common literal modulus and the surviving ray mean}
\label{cbk:ray-proof}
This last construction concerns the pure block only. The primary
section on good ideals is multiplicative because
$\mu_6\to(\mathcal O/3)^\times$ is bijective.
Let $\mathfrak c=(\gamma_c)$ be fixed, $S$-supported and divisible
by $(6)$, large enough for all actual row-ray congruences.
For $\mathfrak q=\mathfrak m\mathfrak n$ coprime to it, put
\[
 \vartheta_{\mathfrak q}(x)=(x/\mathfrak m)_6
                     \overline{(x/\mathfrak n)_6},\qquad
 \tau_{\mathfrak q}=\sum_{x\bmod\mathfrak q}
               \vartheta_{\mathfrak q}(x)e_{\gamma_q}(x).
\]
Every good local factor is nontrivial sextic, so its finite
conductor is exactly $\mathfrak q$, including nonunit frequency
zeros. It need not yet be unit-trivial or an ideal character.
The finite indicator $P_{\mathfrak k^2}$ imposes residual generator
$x\equiv1\pmod3$, original $S$-goodness and the actual ray
condition on $\alpha_{\mathfrak k}^2x$. It is independent of
$\mathfrak m,\mathfrak n$. Literal square-divisor inversion is
\begin{equation}\label{cbk:literal-primal}
 \begin{split}
 \mathcal P_D
 =\sum_{(\mathfrak m,\mathfrak n)=1}c_{\mathfrak m}\bar c_{\mathfrak n}
 \sum_{\substack{\mathfrak k\ \mathrm{good,SF}\\N\mathfrak k\le\sqrt b}}
 \mu(\mathfrak k)\vartheta_{\mathfrak q}(\alpha_{\mathfrak k}^2)
 \sum_{x\in\mathcal O}
 P_{\mathfrak k^2}(x\bmod\mathfrak c)
          \vartheta_{\mathfrak q}(x)G(Nx/L_{\mathfrak k}).
 \end{split}
\end{equation}
The original $(\mathfrak k,\mathfrak q)$ zero is kept.
Auxiliary lattice $x$ need not generate a squarefree ideal.

Use the same positive additive-pairing sign as above. CRT separates
the common-modulus Fourier transform into the $\mathfrak q$ sum
at frequency $\beta\gamma_c^{-1}$ and the $\mathfrak c$ sum at
$\beta\gamma_q^{-1}$. Primitive finite Gauss Fourier proves
\eqref{cbk:ray-Fourier} for \emph{every} $\beta$, including its
zero at $(\beta,\mathfrak q)>1$. This is the finite recombination
of all ray components; it is not a unit root at the larger induced
modulus $\mathfrak c\mathfrak q$.
For $g_L(x)=G(|x|^2/L)$ the Euclidean radial transform is
\[
 \widehat g_L(w)=\pi L\int_0^\infty
                         G(t)J_0(2\pi\sqrt{Lt}|w|)\,dt.
\]
At the common dual vector
$w=2\bar\beta/(\bar\delta\bar\gamma_c\bar\gamma_q)$ this equals
$(\sqrt3/2)L\dot G(LN\beta/(N\mathfrak c Q))$.
The lattice covolume is $(\sqrt3/2)N\mathfrak c Q$.
Consequently, writing $F_{\mathfrak k^2}=P_{\mathfrak k^2}\vartheta_{\mathfrak q}$,
\begin{equation}\label{cbk:literal-Poisson}
 \sum_{x\in\mathcal O}F_{\mathfrak k^2}(x)G(Nx/L)
 =\frac L{N\mathfrak c Q}\sum_{\beta\in\mathcal O}
       \widehat F_{\mathfrak k^2}(\beta)
                            \dot G(LN\beta/(N\mathfrak c Q)).
\end{equation}
No class/root phase is discarded: the original finite function
itself has been transformed. Primitive-conductor descriptions
induced to the common modulus must reproduce precisely this
frequency support, rather than a falsely flat lifted Gauss sum.

Combining the formulas gives the explicit signed literal additive
band
\begin{equation}\label{cbk:literal-band}
 \begin{split}
 \mathcal B_Y={}&\sum_{(\mathfrak m,\mathfrak n)=1}
 c_{\mathfrak m}\bar c_{\mathfrak n}
 \sum_{\substack{\mathfrak k\ \mathrm{good,SF}\\N\mathfrak k\le\sqrt b}}
 \frac{\mu(\mathfrak k)L_{\mathfrak k}\tau_{\mathfrak q}
                    \vartheta_{\mathfrak q}(\gamma_c)}{N\mathfrak c Q}
             \vartheta_{\mathfrak q}(\alpha_{\mathfrak k}^2)\\
 &\quad\cdot\sum_{Y<N\beta\le2Y}
       \overline{\vartheta_{\mathfrak q}(\beta)}
       \widehat P_{\mathfrak k^2}(\beta\gamma_q^{-1})
                   \dot G(L_{\mathfrak k}N\beta/(N\mathfrak c Q)).
 \end{split}
\end{equation}
All finite $\mu(\mathfrak k)$, the original unfactored coefficients
and the full pair/ray/raw masks precede any absolute norm.
There is no gcd triangle or positive completion in this identity.
The frequency band is literal and unit-invariant; it is not simply
the old primitive-ideal band under a new name.

On $\beta=\gamma_c\zeta$, the fixed-residue Fourier frequency is
zero. The factor $\vartheta_{\mathfrak q}(\gamma_c)$ cancels and
\[
 \widehat F_{\mathfrak k^2}(\gamma_c\zeta)
        =r_{\mathfrak k^2}\tau_{\mathfrak q}
                                  \overline{\vartheta_{\mathfrak q}(\zeta)}.
\]
This $S$-local mean is not the global $\beta=0$, which is zero by
primitive nontriviality. A six-unit orbit sum in this radial band
vanishes if $\vartheta_{\mathfrak q}|_{\mu_6}$ is nontrivial.
Otherwise it is six times the sum over \emph{all} integral ideals;
values there are defined with any generator, not a nonexistent
primary generator for an $S$-supported ideal.
At $\varepsilon=1+\omega$, local residue powers give
\[
 \vartheta_{\mathfrak q}(\varepsilon)
       =\varepsilon^{(N\mathfrak m-N\mathfrak n)/6}.
\]
Every good prime-ideal norm is one modulo six, and
$(\prod Np-1)/6\equiv\sum(Np-1)/6\pmod6$, proving
\eqref{cbk:unit-selection}. In the surviving case
$\vartheta_{\mathfrak q}$ is a primitive type-zero ideal character
and $\epsilon_{\mathfrak q}=\tau_{\mathfrak q}/\sqrt Q$.
The exact retained mean channel is
\begin{equation}\label{cbk:literal-mean}
 \begin{split}
 \mathcal B_Y^{S\text{-mean}}
 ={}&\sum_{\substack{(\mathfrak m,\mathfrak n)=1\\
                    N\mathfrak m\equiv N\mathfrak n\ (36)}}
 c_{\mathfrak m}\bar c_{\mathfrak n}
 \sum_{\substack{\mathfrak k\ \mathrm{good,SF}\\N\mathfrak k\le\sqrt b}}
 \mu(\mathfrak k)\vartheta_{\mathfrak q}(\alpha_{\mathfrak k}^2)
 \frac{6r_{\mathfrak k^2}}{N\mathfrak c}
                         \frac{L_{\mathfrak k}}{\sqrt Q}\epsilon_{\mathfrak q}\\
 &\quad\cdot\sum_{Y/N\mathfrak c<N\mathfrak z\le2Y/N\mathfrak c}
       \overline{\vartheta_{\mathfrak q}(\mathfrak z)}
                                      \dot G(L_{\mathfrak k}N\mathfrak z/Q).
 \end{split}
\end{equation}
This includes nonunit and $S$-supported dual modes. The $\mathfrak k$
phase is cubic, $\vartheta_{\mathfrak q}(\mathfrak k)^2$, not one,
and its original cutoff has not been Euler-completed.

For the basic $S$ above $2,3$ and no additional ray condition,
modulo $(6)$ the allowed coordinates $a+b\omega$ are
$(1,0),(1,3),(4,3)$. The fourth primary residue $(4,0)$ is excluded
by the original $2$-adic coprimality condition. Thus $r=3$,
$N\mathfrak c=36$, and $6r/N\mathfrak c=1/2$ exactly. Dropping that
condition would give the incorrect $2/3$. The new exact local
Fourier check has $\sum|\widehat P|^2=36\cdot3=108$; it is a
finite $S$-ring consistency check, not an arithmetic-family or
asymptotic band test.

Proposition~\ref{cbk:equal} controls the whole primal equal-norm
sector. That alone does not control this sector in an individual
literal band or $S$-mean channel. To formulate a literal unequal-norm
problem, impose $N\mathfrak m\ne N\mathfrak n$ on the original pair
sum before applying~\eqref{cbk:literal-primal} and
\eqref{cbk:literal-Poisson}. The projected formula retains unequal
norms congruent modulo $36$ as well as all complementary literal
residues. Neither the full $\mathcal B_Y$ nor its separately projected
unequal-norm version is proved small here. At a pure critical block a
sufficient band estimate would be
\[
 |\mathcal B_Y|\ll_{S,\nu,W,\epsilon}
                         B_W^2 D A_1(DA_1Y)^\epsilon,
\]
with the fixed-profile dependence retained. This is not necessary
for the whole kernel, where bands may cancel. Finite ray/root
recombination alone does not provide it; nor does a nonzero local
mean give a lower bound for the whole band. This identifies the
remaining coherent unequal-norm arithmetic problem without deleting
its means or changing any original conductor, capacity or profile.
